%! TEX root = ../master.tex
\begin{theorem}[Trigonometric integrals]\label{thm:trig-integrals}
  Let $R(x, y)$ be a rational function with no poles on the circle
  $x^2 + y^2 = 1$. Then
  \[
    \int_0^{2\pi} R(\cos\theta, \sin\theta) \dd\theta
    = \oint_{\abs{z} = 1}
      R\left( \frac{z + z^{-1}}{2}, \frac{z - z^{-1}}{2i} \right)
      \frac{\dd z}{iz},
  \]
  and the right-hand side is $2\pi i$ times the sum of the residues of its
  integrand inside $\abs{z} = 1$.
  \thmsource{\cite[ch.~6]{saff-snider}}
\end{theorem}

\begin{proof}
  Parametrise the circle by $z = e^{i\theta}$, $0 \leq \theta \leq 2\pi$. Then
  $\dd z = iz \dd\theta$, $\cos\theta = \frac12(z + z^{-1})$ and
  $\sin\theta = \frac{1}{2i}(z - z^{-1})$, which turns the contour integral
  into the real one. The integrand on the right is rational in $z$, so it has
  finitely many poles, and \cref{thm:residue-theorem} applies.
\end{proof}
